% Set from report/draft.md section 3. The prose is the maintainer's; this file is
% its transit into LaTeX and nothing in it was rewritten.
%
% Every caption states what was measured, on what, and cites the ledger entries
% the numbers come from.
\section{Overview}

\subsection{Architecture}

\begin{figure}[htbp]
  \centering
  \includegraphics[width=\textwidth]{architecture.pdf}
  \caption{The six-layer framework, with the three components that sit outside
  it and the three that are named but not built. \emph{Measured} means the
  ledger contains a measurement of the layer's behaviour, and every such box
  names the foreclosure or measurement entry it rests on; \emph{built} means
  implemented and tested with no measurement bearing on the question the layer
  exists to answer; \emph{hypothetical} means recorded as an open decision and
  not built. A conceptual diagram: it is drawn from the structure of
  \texttt{docs/state.yaml} and \texttt{docs/inherited-measurements.yaml} rather
  than from a numeric result, and the script asserts that all seventeen cited
  ids exist in one of them before it draws. The layer numbering
  \texttt{L1}--\texttt{L6} is Chat-surface vocabulary and is not defined in the
  repository (\texttt{fc-011}).}
  \label{fig:architecture}
\end{figure}

The framework is six layers, and a pass through the loop is the clearest way to
describe them.

A fixation is a gaze direction and a vergence. \textbf{L1} turns it into two eye
rotations under Listing's law, which places the cameras. \textbf{L2} encodes the
resulting pair; \textbf{L3} matches, producing disparity with a variance and a
validity mask. \textbf{L4} scales disparity to metric depth by expanding about
the fixated disparity --- first-order-exact at the fovea, degrading as the
square of the offset. \textbf{L5} decides where attention is allocated;
\textbf{L6} picks the next fixation from the belief's uncertainty. The belief
updates, and the loop repeats.

The layers are not equally weighted in practice. L1 is substantial and was
expensive to get right; L2 and L3 are deliberately minimal; L4 carries more than
its size suggests, through the Jacobian that converts disparity variance into
depth variance; L5 was decided as attentional allocation with the mechanical
branch left untested; L6 is a two-line argmax that, as \S4.2 reports, turned out
not to matter.

Three components sit outside the six-layer stack, and each exists because the
stack alone was insufficient.

\textbf{The sampling model} maps a sample index to a unit ray. It is what makes
the projection swappable: a spherical sensor satisfied the same interface as a
pinhole one without a signature change, because the frame is identical and only
the index-to-direction map differs --- which is the whole of what a projection
is.

\textbf{The belief} is head-anchored and stores directions rather than pixels,
so it survives a change of projection untouched. Its failure mode was
instructive: tied initially to the rectifying rotation, which ignores azimuth by
construction, it silently fused measurements of different world directions into
the same cell, and no test caught it because every control varied elevation.

\textbf{The rig/sensor split.} The original rig type carried baseline and focal
length together, which coupled two thirds of the geometry tests to an intrinsic
they did not need. Separating stereo geometry from the sensor is what made the
rotational content portable.

Listing's tilt is a parameter, not a constant: the default is 0.25 per radian of
vergence, while 0.5 is the value at which vertical disparity in the plane of
regard vanishes exactly. The two are distinct and the distinction is pinned by
test.

\subsection{The experiments}

Thirteen experiments, each pre-registered with a stated falsifier before it was
run. Each declared not only what result would falsify its hypothesis, but what
result would mean the experiment had asked the wrong question --- and that
second clause fired more than once.

\begin{table}[htbp]
  \centering
  \small
  \begin{tabular}{@{}llp{0.50\textwidth}@{}}
    \toprule
    & question & outcome \\
    \midrule
    \textbf{exp001} & Does the gaze objective matter?
      & No --- two objectives, zero agreement, indistinguishable results \\
    \textbf{exp002} & Does variance beat the validity mask?
      & Only marginally; the mask carries the effect \\
    \textbf{exp003} & Does verging to acquire help?
      & No --- worse, with a diagnosed mechanism \\
    \textbf{exp004} & Does the rendered scene match the fixture?
      & Away from discontinuities yes; at them no, and the fixture is the worse
        of the two \\
    \textbf{exp005} & Were the earlier results artefacts of the fixture?
      & The artefact carried the effects rather than masking them \\
    \textbf{exp007} & Do the policy results hold on rendered data?
      & Yes, and variance recovers more strongly than on the fixture \\
    \textbf{exp008} & How much occlusion makes the question visible?
      & Sharpest near a tenth; indiscriminate above a sixth \\
    \textbf{exp009} & What does the epipolar violation cost?
      & More than rectifying does, at 9.4\texttimes{} the matcher time \\
    \textbf{exp010} & Does closing the loop help?
      & Yes, at 41\texttimes{} the cost \\
    \textbf{exp011} & Does it hold at more seeds and more levels?
      & Direction yes, margin no; coverage is the result \\
    \textbf{exp012} & What did the inherited statistical bar test?
      & Not what it was read as testing \\
    \textbf{exp013} & Does foveal weighting help on a complete sphere?
      & No --- strictly lossy, on every seed \\
    \textbf{exp014} & Is the gain allocation or scaling?
      & Re-acquisition, away from discontinuities \\
    \bottomrule
  \end{tabular}
\end{table}

Four of the thirteen --- exp004, exp005, exp008 and exp012 --- are about the
instrument or the criterion rather than about the framework. They are not a
detour. Each changed how the experiments around it had to be read, and two of
them reversed conclusions that had already been drawn.

\subsection{The instrument}

\textbf{This section reports a result, not background.} The most surprising
finding in the project concerns the stimulus rather than the framework, and
every result after it had to be re-read in its light.

\begin{figure}[htbp]
  \centering
  \includegraphics[width=\textwidth]{fixture.pdf}
  \caption{The analytic fixture at seed 0: the stereo pair and the ground-truth
  depth field behind it, at the parameters every analytic-fixture result in this
  report was measured at --- $240\times320$, $f = 700$\,px, baseline 65\,mm.
  Depth is present only in disparity; neither monocular image contains it, which
  is the property Julesz's construction was designed
  for~\cite{julesz1960binocular}. \emph{Regenerated} from
  \texttt{make\_synthetic\_scene(seed=0)}, ported unchanged from
  \texttt{visgraf/bioeye@e908170}~\cite{bioeye}. The figure asserts
  \texttt{gap-010}'s first mechanism before drawing: zero right-image values
  fall outside the left image's range, because the right image is resampled from
  the left rather than rendered.}
  \label{fig:fixture}
\end{figure}

Random-dot stereograms are the right synthetic fixture for this work, for the
reason Julesz introduced them: depth exists only in disparity, so nothing can be
recovered by exploiting shading or texture gradient. The fixture used here
follows the standard construction --- generate a texture, then warp it by the
ground-truth disparity field to produce the second view.

That construction has a defect, and it is not the obvious one.

At a depth discontinuity the disparity field jumps, so adjacent output pixels
draw from widely separated input columns. The second image acquires a stretched
or duplicated patch of real texture with no consistent disparity. The matcher
finds it. The result is a confident, incorrect depth estimate --- and it
survives the front end's validity test, of which left--right consistency is one
of two conditions: 78.5\% of discontinuity pixels are marked valid, and 21.7\%
of those are wrong by more than two pixels.

The obvious criticism of a warped stereogram is that it lacks true
half-occlusions. The measured problem is worse and has the opposite sign: the
fixture substitutes an artefact that is harder than real occlusion and invisible
to the detector meant to catch it. Rendering the same geometry in a
physically-based renderer, where occluded pixels genuinely have no correspondent
and fail consistency, the fixture is 23\texttimes{} worse in the tail at depth
discontinuities and 1.6\texttimes{} worse in the median. The damage is
concentrated in a small number of badly wrong estimates rather than spread
across the field --- which is what an artefact that passes validity looks like.

The pixels within ten of a depth discontinuity are 38\% of the valid measured
pixels and carry 82\% of the total squared error.

\begin{figure}[htbp]
  \centering
  \includegraphics[width=\textwidth]{diagnostic.pdf}
  \caption{The same geometry from two sources, at seed 0, under the ported loop
  with policy \texttt{A'} for 18 fixations: the analytic fixture (F), whose
  right image is resampled from its left, and a Blender render (R), where
  occluded pixels genuinely have no correspondent. The fixture's error
  concentrates in bands at the depth discontinuities that the render's does not
  have, and its RMSE settles at 0.475\,m against the render's 0.225\,m.
  \emph{Measured}: exp004, whose recorded band figures are \texttt{AT} p90
  1.4982\,m for F against 0.0734\,m for R
  (\texttt{experiments/exp004\_scene\_model\_check/results.json}). The panels are
  regenerated from committed code at the pinned seed because exp004 stored no
  arrays; both regenerated scanpaths are asserted against the two the experiment
  recorded before the figure is drawn, and the two loops agree on 1 of 18
  fixations.}
  \label{fig:diagnostic}
\end{figure}
